\section{Central sequence algebras}\label{sec:central}
We now transfer the full signs of Theorem~\ref{thm:tensor-intro} to the
central sequence algebra and apply the absorption criterion of Robert
and Tikuisis. This completes the implication from finite nuclear dimension
to absorption.

Retain the free ultrafilter \(\omega\) on \(\N\). For a \(C^*\)-algebra \(A\), let
\[
 A_\omega=\ell^\infty(A)/\{(a_l):\lim_{l\to\omega}\|a_l\|=0\},\qquad
 F(A)=(A_\omega\cap A')/\ann(A),
\]
where \(A\) is embedded diagonally and
\[
 \ann(A)=\{b\in A_\omega:ba=ab=0\text{ for every }a\in A\}.
\]
For nonzero separable \(A\), this quotient is nonzero and unital. If \((e_l)\) is a sequential positive contractive approximate identity, its class is central and represents the unit. Indeed, for a bounded central sequence \((b_l)\) and \(a\in A\), commuting \(b_l\) past \(a\) shows \((e_lb_l-b_l)a\to_\omega0\), and likewise on the other side. Also \(e_la\to a\), so the unit is nonzero.

We use the following two results.
\begin{theorem}[Winter--Zacharias and Robert--Tikuisis]\label{thm:imports}
Let \(A\ne0\) be separable with \(\nucdim A\leq d<\infty\).
\begin{enumerate}
\item If no hereditary subalgebra of \(A\) has a finite-dimensional
irreducible representation, then there are finitely many finite-dimensional
algebras $E_j$, completely positive order-zero maps
$\varphi_j:E_j\to F(A)$, and elements $c_j\in F(A)$ such that every
matrix summand of every $E_j$ has size at least two and
\[
 1_{F(A)}=\sum_jc_j^*\varphi_j(1_{E_j})c_j.
\]
This is the matrix-threshold-two consequence of the covering-number
bound $\operatorname{cov}(F(A))\le(d+1)^2$
\cite[Definition~4.1 and Proposition~4.3]{WZNuclear}.
\item If \(F(A)\) contains two full orthogonal positive elements, then \(A\cong A\otimes_{\min}\Z\) \cite[Theorem 1.2]{RT}.
\end{enumerate}
Both assertions apply to nonunital \(A\), with the annihilator quotient just defined.
\end{theorem}

\begin{lemma}\label{lem:nocharacters}
If \(A\ne0\) is separable, nowhere scattered, and has finite nuclear dimension, then \(F(A)\) has no characters.
\end{lemma}
\begin{proof}
A hereditary subalgebra \(H\) is strongly Morita equivalent to the ideal \(I\) it generates. If \(H\) had a matrix quotient, ideal correspondence would give an ideal \(J\subseteq I\) for which \(I/J\) is strongly Morita equivalent to a matrix algebra, hence to \(\C\). By imprimitivity, \(I/J\) would be elementary. Ideals of \(I\) are ideals of \(A\), contradicting nowhere scatteredness. This is also the equivalence in \cite[Theorem 3.1]{TV}.

Apply Theorem~\ref{thm:imports}(1) to $A$, put $B=F(A)$, and retain its
maps $\varphi_j:E_j\to B$, elements $c_j$, and displayed identity.

If \(\chi:B\to\C\) were a character, its compositions with these maps would be order-zero positive functionals. Such a functional on \(M_r\), \(r\geq2\), is zero. To see this, a nonzero positive functional has a positive density matrix with an eigenvector \(\xi\) of positive eigenvalue. For a unit eigenvector \(\eta\perp\xi\), the orthogonal rank-one projections onto \((\xi+\eta)/\sqrt2\) and \((\xi-\eta)/\sqrt2\) both have strictly positive values, contrary to order zero. Thus \(\chi\circ\varphi_j=0\) for every \(j\), contradicting the displayed identity and \(\chi(1_B)=1\).
\end{proof}

\begin{lemma}[Commuting reindexed copies]\label{lem:reindex}
Let \(A\ne0\) be separable, and let \(D\subseteq F(A)\) be a separable unital subalgebra with the same unit. For every finite \(n\), there is a unital \(*\)-homomorphism
\[
 D^{\otimes_{\max}n}\longrightarrow F(A).
\]
\end{lemma}
\begin{proof}
It suffices to prove that for every separable subalgebra \(C\subseteq F(A)\), there is a unital \(*\)-homomorphism from \(D\) into \(F(A)\) with range commuting with \(C\). This is Kirchberg's reindexing mechanism, in the form recorded in
\cite[Proposition~2.2 and Corollary~2.3]{KR} for unital $A$. We include
the annihilator details needed for nonunital $A$.

Choose a countable dense unital \(*\)-algebra \(D_0\subset D\) over \(\Q+i\Q\). For each \(x\in D_0\), choose a bounded lift \((a_l^x)_l\) in \(A\) that represents a central element mapping to \(x\), with
\begin{equation}\label{eq:liftbound}
 \sup_l\|a_l^x\|\leq2\|x\|.
\end{equation}
Take the zero lift for zero and the approximate identity for \(1\). The quotient norm permits these choices: first lift within a bound strictly below \(2\|x\|\), then truncate coordinate norms without changing the ultrapower class. The additive, scalar, multiplicative, and adjoint relation defects between these lifts annihilate \(A\) on both sides along \(\omega\).

Choose a countable dense family in \(C\) with bounded lifts \((b_s^u)_s\) in \(A\). Also fix a countable dense subset of \(A\). List the following countably many requirements on an old index \(l\): centrality of each \(a_l^x\) against that subset; left and right products of every relation defect with that subset tending to zero; and
\[
 a_l^1a\longrightarrow a,\qquad aa_l^1\longrightarrow a
\]
for its elements \(a\). All these requirements hold along \(\omega\).

For each output index \(s\), choose \(l(s)\) so that the first \(s\) requirements hold with error less than \(1/s\), and also
\[
 \|[a_{l(s)}^x,b_s^u]\|<1/s
\]
for the first \(s\) pairs \((x,u)\) in an enumeration. For this fixed \(s\), the finitely many \(b_s^u\) are fixed elements of \(A\). Original centrality therefore makes all these commutator conditions \(\omega\)-large as conditions on \(l\). A finite intersection with the other \(\omega\)-large conditions is nonempty, so the choice is possible.

The new sequences \((a_{l(s)}^x)_s\) are central, and all their relation defects annihilate \(A\) on both sides: the tests hold in the ordinary limit, and the uniform bounds extend them from the dense subset to all of \(A\). Thus they define a unital rational-complex \(*\)-homomorphism on \(D_0\) into \(F(A)\). The unit assertion follows from the approximate-identity tests. Bound \eqref{eq:liftbound} permits continuous extension to a complex unital \(*\)-homomorphism on \(D\). Its range commutes with the dense family in \(C\), hence with \(C\). Injectivity is not needed.

Starting with one copy and repeatedly taking \(C\) to be generated by the previous images gives \(n\) commuting unital images of \(D\). The universal property of the maximal tensor product gives the stated homomorphism.
\end{proof}

\begin{proof}[Proof of Theorem~\ref{thm:main}]
The assertion is immediate for \(A=0\), so assume \(A\ne0\).
Assume first that $\nucdim(A)<\infty$; we prove (1)$\Rightarrow$(2).
By Lemma~\ref{lem:nocharacters}, \(B=F(A)\) has no characters. Choose in \(B\) a finite commutator witness as in \eqref{eq:commutator}, and let \(D\) be the unital \(C^*\)-subalgebra generated by all its entries. Then \(D\) is separable and has no characters, since applying a character of \(D\) to that inequality would give \(1<1/2\).

By Theorem~\ref{thm:tensor-intro}, some \(D^{\otimes_{\max}n}\) contains full orthogonal positive elements \(a,b\). Lemma~\ref{lem:reindex} supplies a unital homomorphism into \(B\). Their images remain orthogonal and full: applying a unital homomorphism to finite ideal combinations approximating the unit proves that fullness passes to the target. Theorem~\ref{thm:imports}(2) now gives \(A\cong A\otimes_{\min}\Z\).
Conversely, for separable nuclear $A$, absorption and
Theorem~\ref{intro:main} imply $\nucdim(A)\le1$. This bound implies
finite nuclear dimension, completing all three equivalences in
Theorem~\ref{thm:main}.
\end{proof}

\begin{remark}
The same argument proves that \(F(A)\) has two full orthogonal positive elements whenever \(A\ne0\) is separable and \(F(A)\) has no characters. Finite nuclear dimension enters the proof of Theorem~\ref{thm:main} only through the two imported results in Theorem~\ref{thm:imports}.
\end{remark}

\subsection{Cuntz-semigroup consequence}
\label{subsec:cu-consequence}

We use the completed stabilized Cuntz semigroup $\Cu(X)$ formed from
$(X\otimes\K)_+$, with no distinguished scale. Isomorphisms in
$\mathbf{Cu}$ preserve zero, addition, order, increasing sequential
suprema, and the way-below relation in both directions. The isomorphism
in Corollary~\ref{cor:regularity-cu}\textup{(ii)} need not agree with
$\Cu(\iota_A)$.

\begin{proof}[Proof of Corollary~\ref{cor:regularity-cu}]
Theorem~\ref{thm:main} gives the equivalence of its three conditions.
The companion \cite[Corollaries~1.3 and~1.4]{CuntzCompanion} gives
that \textup{(i)} implies absorption and that \textup{(ii)} is equivalent
to absorption, even without the exclusion on elementary ideal subquotients.

\Needspace{7\baselineskip}
To prove that absorption implies \textup{(i)}, the case $A=0$ is immediate.
Put $D=A\otimes\Z$ and $j_X(x)=x\otimes1_\Z$, so $j_A=\iota_A$.
Strong self-absorption gives an isomorphism
$\varphi:\Z\to\Z\otimes\Z$ approximately unitarily equivalent to
$z\mapsto z\otimes1_\Z$. The companion's nonunital calculation
\cite[proof of Corollary~1.4]{CuntzCompanion} gives
$\Cu(j_D)=\Cu(\id_A\otimes\varphi)$, an isomorphism: the implementing
unitaries $u_n\in\Z\otimes\Z$ act through
$1_{M(A)}\otimes u_n\in M(A\otimes\Z\otimes\Z)$, and positive cuts
and approximate units after stabilization give equality of the completed
Cuntz maps.

For an absorption isomorphism $\theta:A\to D$, both sides of the first
identity below send $a$ to $\theta(a)\otimes1_\Z$. Functoriality then gives
\[
\begin{aligned}
 (\theta\otimes\id_\Z)\circ j_A&=j_D\circ\theta,\\
 \Cu(j_A)&=\Cu(\theta\otimes\id_\Z)^{-1}
       \circ\Cu(j_D)\circ\Cu(\theta).
\end{aligned}
\]
All three factors on the right are isomorphisms, proving \textup{(i)}.
Thus \textup{(ii)} reaches the canonical condition \textup{(i)} through
absorption, not by identifying its abstract isomorphism with $\Cu(\iota_A)$.
\end{proof}
